% Perda de validação do ajuste supervisionado. Três pontos apenas — uma época é
% a unidade de avaliação —, mas é a série que decide qual checkpoint vira o
% modelo entregue, e o mínimo na época 2 é visível de imediato.
\begin{figure}[htb]
	\Caption{\label{fig:sft-eval} Perda de validação por época no ajuste supervisionado}
	\centering
	\begin{tikzpicture}
		\begin{axis}[
			width=0.62\textwidth, height=5.2cm,
			xlabel={Época}, ylabel={Perda de validação},
			xmin=0.7, xmax=3.3, ymin=0.408, ymax=0.448,
			xtick={1,2,3},
			ytick={0.41,0.42,0.43,0.44},
			yticklabel style={/pgf/number format/fixed, /pgf/number format/precision=3},
			ymajorgrids, grid style={draw=black!12},
			axis lines*=left, tick align=outside,
			tick label style={font=\small}, label style={font=\small},
			nodes near coords, nodes near coords style={font=\scriptsize, /pgf/number format/fixed, /pgf/number format/precision=4},
			every node near coord/.append style={anchor=south west, xshift=2pt},
		]
			\addplot[corSFT, thick, mark=*, mark size=2.6pt]
				table[x=epoca, y=eval_loss, col sep=comma] {figuras/dados/sft_eval.csv};
			\draw[corTeto, dashed] (axis cs:2,0.408) -- (axis cs:2,0.416);
			\node[corTeto, font=\scriptsize, anchor=north] at (axis cs:2,0.4128) {\textit{checkpoint} entregue};
		\end{axis}
	\end{tikzpicture}
	\Fonte{Elaborado pelo autor}
	\Nota{Validação de 78 exemplos, agrupada por documento de origem. O mínimo na
	época 2 determinou o adaptador entregue: a subida na época 3 indica início de
	sobreajuste, e o critério de parada por validação descartou-a
	automaticamente. Verificado por soma de verificação que o adaptador salvo é
	idêntico ao \textit{checkpoint} da época 2.}
\end{figure}
